% ============================================================
%  Remote Healthcare Monitoring System
%  Graduation Project Report
%  University of Baghdad — Al-Khwarizmi College of Engineering
%  Mechatronics Engineering
% ============================================================
\documentclass[12pt,a4paper]{report}

% ── Packages ────────────────────────────────────────────────
\usepackage[utf8]{inputenc}
\usepackage[T1]{fontenc}
\usepackage[english]{babel}
\usepackage{geometry}
\usepackage{graphicx}
\usepackage{float}
\usepackage{booktabs}
\usepackage{array}
\usepackage{longtable}
\usepackage{hyperref}
\usepackage{fancyhdr}
\usepackage{setspace}
\usepackage{titlesec}
\usepackage{caption}
\usepackage{subcaption}
\usepackage{enumitem}
\usepackage{amsmath}
\usepackage{xcolor}
\usepackage{tikz}
\usetikzlibrary{shapes,arrows.meta,positioning,fit}

\geometry{left=2.5cm, right=2.5cm, top=2.5cm, bottom=2.5cm}
\onehalfspacing
\hypersetup{colorlinks=true, linkcolor=blue, citecolor=blue, urlcolor=blue}

% ── Header / Footer ─────────────────────────────────────────
\pagestyle{fancy}
\fancyhf{}
\fancyhead[L]{\small Remote Healthcare Monitoring System}
\fancyhead[R]{\small Mechatronics Engineering}
\fancyfoot[C]{\thepage}
\renewcommand{\headrulewidth}{0.4pt}

% ── Helper: insert figure or placeholder ────────────────────
% Usage: \projectfigure{filename}{height}{caption}{label}
\newcommand{\projectfigure}[4]{%
  \begin{figure}[H]
    \centering
    \IfFileExists{figures/#1}{%
      \includegraphics[height=#2,keepaspectratio]{figures/#1}%
    }{%
      \fcolorbox{gray!60}{gray!10}{%
        \parbox[c][#2][c]{0.88\textwidth}{%
          \centering\small
          \textbf{[INSERT IMAGE]}\\[4pt]
          \texttt{figures/#1}\\[6pt]
          See \texttt{PHOTOS\_GUIDE.md} for details.
        }%
      }%
    }
    \caption{#3}
    \label{#4}
  \end{figure}%
}

% ── Title page info (edit as needed) ────────────────────────
\newcommand{\ProjectTitle}{Remote Healthcare Monitoring System}
\newcommand{\Department}{Mechatronics Engineering}
\newcommand{\College}{Al-Khwarizmi College of Engineering}
\newcommand{\University}{University of Baghdad}
\newcommand{\Ministry}{Ministry of Higher Education and Scientific Research}
\newcommand{\ReportYear}{2025--2026}

\newcommand{\SupervisorOne}{Dr. Ahmed Rahman}
\newcommand{\SupervisorTwo}{Dr. Adnan Jabar}
\newcommand{\StudentOne}{Mahmood Hamid}
\newcommand{\StudentTwo}{Sinan Raid}
\newcommand{\StudentThree}{Dhuha Kareem}

% ============================================================
\begin{document}

% ── TITLE PAGE ──────────────────────────────────────────────
\begin{titlepage}
  \centering
  \vspace*{0.5cm}

  \IfFileExists{figures/university-logo.png}{%
    \includegraphics[height=2.2cm]{figures/university-logo.png}\hspace{1.5cm}%
  }{%
    \fbox{\parbox[c][2.2cm][c]{3cm}{\centering\small University\\Logo}}%
  }%
  \IfFileExists{figures/college-logo.png}{%
    \includegraphics[height=2.2cm]{figures/college-logo.png}%
  }{%
    \fbox{\parbox[c][2.2cm][c]{3cm}{\centering\small College\\Logo}}%
  }

  \vspace{1cm}
  {\Large \Ministry\par}
  \vspace{0.4cm}
  {\LARGE\bfseries \University\par}
  \vspace{0.3cm}
  {\Large \College\par}
  \vspace{0.3cm}
  {\large \Department\par}

  \vspace{2cm}
  {\Huge\bfseries \ProjectTitle\par}

  \vspace{0.8cm}
  {\Large Graduation Project Report\par}

  \vspace{2cm}
  \begin{minipage}{0.45\textwidth}
    \centering
    \textbf{Supervisors}\\[6pt]
    \SupervisorOne\\[4pt]
    \SupervisorTwo
  \end{minipage}\hfill
  \begin{minipage}{0.45\textwidth}
    \centering
    \textbf{Submitted By}\\[6pt]
    \StudentOne\\[4pt]
    \StudentTwo\\[4pt]
    \StudentThree
  \end{minipage}

  \vfill
  {\large Baghdad --- \ReportYear\par}
\end{titlepage}

% ── FRONT MATTER ────────────────────────────────────────────
\pagenumbering{roman}
\setcounter{page}{1}

\chapter*{Abstract}
\addcontentsline{toc}{chapter}{Abstract}

Healthcare systems worldwide are shifting toward continuous, technology-assisted patient monitoring outside traditional hospital settings. This graduation project presents a \textbf{Remote Healthcare Monitoring System (RHMS)}---an integrated platform that connects wearable smartwatch devices, mobile applications, and a web-based dashboard through a secure cloud backend.

The system uses a Samsung Galaxy Watch running Wear OS to collect vital signs including heart rate, blood oxygen saturation (SpO$_2$), skin/body temperature, step count, and fall events. Sensor data is transmitted in real time via the MQTT protocol to an ASP.NET Core backend, stored in PostgreSQL, and pushed to connected clients through SignalR. When readings exceed configurable thresholds, the system automatically generates alerts and sends push notifications to assigned doctors and family members.

The platform supports four user roles---Admin, Doctor, Patient, and Relative---enabling multi-stakeholder remote care. Results demonstrate successful end-to-end data flow from wearable sensors to live dashboards, reliable alert triggering, and a scalable Docker-based deployment architecture.

\vspace{0.5cm}
\noindent\textbf{Keywords:} Remote Patient Monitoring, Wearable Technology, IoT, MQTT, SignalR, Vital Signs, Fall Detection, Telemedicine, Smartwatch.

\chapter*{Acknowledgments}
\addcontentsline{toc}{chapter}{Acknowledgments}

We would like to express our sincere gratitude to our supervisors, \SupervisorOne\ and \SupervisorTwo, for their guidance, support, and valuable feedback throughout this project. We also thank the staff of \College\ for providing the academic environment and resources that made this work possible. Finally, we thank our families and colleagues for their encouragement during the development of this system.

\tableofcontents
\listoffigures
\listoftables

% ── MAIN BODY ───────────────────────────────────────────────
\clearpage
\pagenumbering{arabic}

% ============================================================
\chapter{Introduction}
% ============================================================

\section{Background}

Modern healthcare faces growing pressure from an aging population, rising rates of chronic disease, and limited hospital capacity. Traditional care models rely on periodic clinic visits and manual vital-sign measurement, which cannot detect sudden deterioration between appointments. Remote Patient Monitoring (RPM) addresses this gap by using connected devices to track health metrics continuously and alert caregivers when intervention is needed \cite{rpm_survey}.

Wearable smartwatches equipped with optical heart-rate sensors, pulse oximeters, accelerometers, and temperature sensors have become practical tools for consumer health tracking. Combined with mobile networks and cloud computing, these devices enable real-time transmission of physiological data to healthcare providers and family members anywhere in the world.

\section{Problem Statement}

Patients with cardiovascular conditions, respiratory illness, diabetes, or mobility limitations require ongoing supervision. Manual monitoring is labor-intensive, expensive, and often delayed. Critical events such as hypoxemia, tachycardia, hyperthermia, or falls may go unnoticed for hours if no caregiver is present. There is a clear need for an automated, connected system that:

\begin{itemize}[leftmargin=*]
  \item Continuously measures key vital signs using wearable hardware.
  \item Transmits data securely and reliably to a central server.
  \item Presents live readings to doctors, patients, and relatives.
  \item Triggers immediate alerts when abnormal conditions are detected.
\end{itemize}

\section{Project Objectives}

The main objectives of this graduation project are:

\begin{enumerate}[leftmargin=*]
  \item Design and implement a complete remote healthcare monitoring system using wearable and software technologies.
  \item Integrate Samsung Galaxy Watch sensors (heart rate, SpO$_2$, temperature) with a Wear OS application.
  \item Develop a scalable backend API for data ingestion, storage, alert evaluation, and user management.
  \item Build web and mobile client applications for real-time vital monitoring and alert notification.
  \item Implement fall detection using accelerometer data from the smartwatch.
  \item Deploy the system using containerized infrastructure (Docker) for reproducibility.
  \item Test and validate end-to-end functionality under realistic usage scenarios.
\end{enumerate}

\section{Scope and Limitations}

This project focuses on prototype-level RPM for demonstration and academic evaluation. Blood pressure and ECG readings are supported in the data model but depend on hardware availability. The system is not certified as a medical device and is intended for research and educational purposes. Production deployment would require regulatory compliance, clinical validation, and enhanced security hardening.

\projectfigure{project-overview.png}{7cm}{Complete project setup: smartwatch, mobile app, and web dashboard running simultaneously}{fig:overview}

\section{Report Organization}

Chapter~\ref{ch:litreview} reviews related work in remote monitoring and IoT healthcare. Chapter~\ref{ch:requirements} defines system requirements. Chapter~\ref{ch:design} presents the architecture and design. Chapter~\ref{ch:implementation} describes implementation details. Chapter~\ref{ch:testing} covers testing and results. Chapter~\ref{ch:conclusion} concludes with future work.

% ============================================================
\chapter{Literature Review}
\label{ch:litreview}
% ============================================================

\section{Remote Patient Monitoring}

Remote Patient Monitoring has been widely studied as a means to reduce hospital readmissions and improve chronic disease management. Research shows that continuous monitoring of heart rate, blood pressure, and oxygen saturation can detect early signs of deterioration in patients with heart failure, COPD, and post-surgical recovery \cite{rpm_chronic}.

Telemedicine platforms typically combine wearable sensors, mobile applications, and cloud dashboards. Key challenges include data accuracy, network reliability, patient adherence (device wearing), and integration with existing healthcare workflows.

\section{Wearable Health Technology}

Consumer and medical-grade wearables now support photoplethysmography (PPG) for heart rate and SpO$_2$, skin temperature sensing, and activity tracking. Samsung's Health Sensor SDK exposes low-level sensor access on Galaxy Watch devices, enabling third-party applications to stream clinical-style measurements in real time.

Accelerometer-based fall detection algorithms analyze sudden impact followed by post-impact immobility or low acceleration (free-fall signature). These methods are widely used in elderly care and assistive monitoring systems \cite{fall_detection}.

\section{IoT Communication Protocols}

MQTT (Message Queuing Telemetry Transport) is a lightweight publish/subscribe protocol designed for IoT devices with limited bandwidth. Its low overhead and Quality of Service (QoS) levels make it suitable for streaming vital-sign payloads from wearables to cloud servers \cite{mqtt_iot}.

For browser and mobile clients requiring bidirectional real-time updates, WebSocket-based technologies such as SignalR provide efficient server-push capabilities, complementing MQTT's device-to-server communication.

\section{Related Systems}

Commercial RPM solutions (e.g., Philips eCareCoordinator, BioIntelliSense) offer enterprise-grade monitoring but are costly and closed-source. Open-source and academic projects demonstrate similar architectures: edge device $\rightarrow$ message broker $\rightarrow$ REST API $\rightarrow$ dashboard. This project follows that proven pattern while using modern .NET, React, and Kotlin stacks.

% ============================================================
\chapter{System Analysis and Requirements}
\label{ch:requirements}
% ============================================================

\section{Stakeholders and User Roles}

The system serves four distinct user roles:

\begin{table}[H]
  \centering
  \caption{User roles and primary responsibilities}
  \label{tab:roles}
  \begin{tabular}{@{}llp{7cm}@{}}
    \toprule
    \textbf{Role} & \textbf{Actor} & \textbf{Responsibilities} \\
    \midrule
    Admin       & System administrator & Manage users, devices, and system configuration \\
    Doctor      & Healthcare provider    & Monitor assigned patients, review alerts, set thresholds \\
    Patient     & Monitored individual   & Wear device, view own vitals, receive alerts \\
    Relative    & Family member          & View linked patient's vitals and receive alerts \\
    \bottomrule
  \end{tabular}
\end{table}

\section{Functional Requirements}

\begin{enumerate}[label=FR\arabic*:, leftmargin=*]
  \item The watch application shall measure heart rate, SpO$_2$, and skin/body temperature.
  \item The watch application shall detect falls using accelerometer data when the device is worn.
  \item The watch application shall publish vital readings to the server via MQTT.
  \item The backend shall ingest, validate, and persist all incoming vital records.
  \item The backend shall evaluate readings against per-patient alert thresholds.
  \item The system shall generate alerts for abnormal heart rate, low SpO$_2$, high temperature, high blood pressure, and fall events.
  \item The system shall push real-time vital updates to web and mobile clients via SignalR.
  \item The system shall send push notifications (FCM) to doctors and relatives on critical alerts.
  \item Users shall register, authenticate, and access features according to their role (JWT).
  \item Doctors shall be assigned to patients; relatives shall be linked to patients.
  \item The web dashboard shall display live and historical vital data.
  \item Administrators shall manage user accounts and system settings.
\end{enumerate}

\section{Non-Functional Requirements}

\begin{enumerate}[label=NFR\arabic*:, leftmargin=*]
  \item \textbf{Performance:} Vital updates shall reach dashboards within 2 seconds under normal network conditions.
  \item \textbf{Scalability:} Backend shall support horizontal scaling via Redis SignalR backplane.
  \item \textbf{Security:} All API access shall require JWT authentication; passwords shall be hashed.
  \item \textbf{Availability:} Docker-based deployment with health checks for all services.
  \item \textbf{Maintainability:} Clean Architecture separation (Domain, Application, Infrastructure, API).
  \item \textbf{Usability:} Interfaces shall be readable on watch, phone, and desktop screens.
\end{enumerate}

\section{Vital Signs and Alert Thresholds}

Table~\ref{tab:vitals} lists monitored metrics. Default alert thresholds are configurable per patient.

\begin{table}[H]
  \centering
  \caption{Monitored vital signs and default alert thresholds}
  \label{tab:vitals}
  \begin{tabular}{@{}llll@{}}
    \toprule
    \textbf{Metric} & \textbf{Unit} & \textbf{Source} & \textbf{Default Threshold} \\
    \midrule
    Heart Rate      & bpm   & PPG sensor       & 40--130 bpm \\
    SpO$_2$         & \%    & Pulse oximeter   & Min 90\% \\
    Body Temperature& \textdegree C & Skin sensor & Max 38.5\textdegree C \\
    Systolic BP     & mmHg  & (optional)       & Max 160 mmHg \\
    Fall Detection  & bool  & Accelerometer    & Any fall event \\
    Steps / Calories& count & Step counter     & Activity tracking \\
    \bottomrule
  \end{tabular}
\end{table}

% ============================================================
\chapter{System Design}
\label{ch:design}
% ============================================================

\section{Overall Architecture}

The RHMS follows a four-tier IoT architecture:

\begin{enumerate}[leftmargin=*]
  \item \textbf{Perception Layer:} Samsung Galaxy Watch sensors (PPG, SpO$_2$, temperature, accelerometer).
  \item \textbf{Edge Layer:} Wear OS application with foreground service for continuous monitoring and MQTT publishing.
  \item \textbf{Network \& Service Layer:} MQTT broker (Mosquitto), ASP.NET Core API, SignalR hubs, PostgreSQL, Redis.
  \item \textbf{Application Layer:} React web app, Android mobile app for doctors/patients/relatives.
\end{enumerate}

\projectfigure{system-architecture.png}{8cm}{System architecture: wearable device, MQTT broker, backend services, and client applications}{fig:architecture}

\section{Data Flow}

The end-to-end data flow is illustrated in Figure~\ref{fig:dataflow}:

\begin{enumerate}[leftmargin=*]
  \item Patient wears the smartwatch and starts the monitoring service.
  \item \texttt{VitalsMonitorService} reads active sensor (HR, SpO$_2$, or temperature) via Samsung Health Sensor SDK.
  \item \texttt{MotionSensorHub} concurrently tracks steps and evaluates accelerometer data for falls.
  \item A JSON payload is published to MQTT topic \texttt{vitals/\{patientId\}/data}.
  \item Backend \texttt{MqttBackgroundService} subscribes, deserializes, and dispatches \texttt{IngestVitalCommand}.
  \item A \texttt{VitalRecord} entity is persisted; \texttt{VitalRecordedEvent} triggers threshold evaluation.
  \item If alerts are created, \texttt{AlertTriggeredEvent} sends FCM push and SignalR broadcast.
  \item Web and mobile clients receive live updates through the \texttt{VitalsHub} SignalR connection.
\end{enumerate}

\projectfigure{data-flow.png}{7cm}{Data flow from sensor reading to dashboard update and alert notification}{fig:dataflow}

\section{Use Case Model}

Figure~\ref{fig:usecase} shows the primary interactions between actors and the system.

\projectfigure{use-case-diagram.png}{7cm}{UML use-case diagram for Patient, Doctor, Relative, and Admin actors}{fig:usecase}

\section{Backend Design (Clean Architecture)}

The backend is organized into four layers:

\begin{itemize}[leftmargin=*]
  \item \textbf{RPM.Domain:} Entities (\texttt{VitalRecord}, \texttt{Alert}, \texttt{User}), enums, domain events.
  \item \textbf{RPM.Application:} CQRS commands/queries via MediatR, DTOs, business handlers.
  \item \textbf{RPM.Infrastructure:} EF Core persistence, MQTT service, FCM, Redis cache, MinIO storage.
  \item \textbf{RPM.API:} REST controllers, SignalR hubs (\texttt{VitalsHub}, \texttt{ChatHub}), JWT middleware.
\end{itemize}

Domain events (\texttt{VitalRecordedEvent}, \texttt{AlertTriggeredEvent}) decouple vital ingestion from alert processing and notification delivery.

\section{Database Design}

Core entities and relationships:

\begin{itemize}[leftmargin=*]
  \item \texttt{User} $\rightarrow$ role-specific profiles (\texttt{PatientProfile}, \texttt{DoctorProfile}).
  \item \texttt{PatientProfile} $\rightarrow$ many \texttt{VitalRecord}, \texttt{Alert}, \texttt{Device}.
  \item \texttt{VitalRecord} $\rightarrow$ optional linked \texttt{Alert}(s).
  \item \texttt{AlertThreshold} $\rightarrow$ one per patient (configurable min/max values).
  \item \texttt{PatientDoctorAssignment}, \texttt{PatientRelativeLink} $\rightarrow$ care-team relationships.
\end{itemize}

\section{Fall Detection Algorithm}

The watch application implements a two-phase accelerometer algorithm in \texttt{MotionSensorHub}:

\begin{enumerate}[leftmargin=*]
  \item \textbf{Impact detection:} Acceleration magnitude exceeds $\sim$3.5\,g or jerk exceeds threshold.
  \item \textbf{Free-fall confirmation:} After peak impact, sustained low acceleration ($<$1.2\,g) for $>$120\,ms indicates a fall.
  \item Falls are only counted when \texttt{isWearing} is true (watch on wrist).
  \item A fall flag remains active for 15 seconds to ensure server ingestion.
\end{enumerate}

\section{Deployment Architecture}

All services run in Docker containers orchestrated by Docker Compose:

\begin{table}[H]
  \centering
  \caption{Docker infrastructure services}
  \label{tab:docker}
  \begin{tabular}{@{}llp{5.5cm}@{}}
    \toprule
    \textbf{Service} & \textbf{Port} & \textbf{Purpose} \\
    \midrule
    RPM API          & 8080  & REST API + SignalR hubs \\
    PostgreSQL       & 5432  & Primary database (TimescaleDB) \\
    Redis            & 6379  & Cache + SignalR backplane \\
    Mosquitto        & 1883  & MQTT message broker \\
    MinIO            & 9000  & Object storage \\
    Seq              & 5341  & Centralized logging \\
    \bottomrule
  \end{tabular}
\end{table}

% ============================================================
\chapter{Implementation}
\label{ch:implementation}
% ============================================================

\section{Technology Stack}

\begin{table}[H]
  \centering
  \caption{Technologies used in each component}
  \label{tab:tech}
  \begin{tabular}{@{}ll@{}}
    \toprule
    \textbf{Component} & \textbf{Technologies} \\
    \midrule
    Smartwatch App     & Kotlin, Wear OS, Hilt DI, Samsung Health Sensor SDK, MQTT \\
    Android Mobile App & Kotlin, Jetpack Compose, SignalR, Retrofit, FCM \\
    Web Application    & React, Vite, SignalR JavaScript client \\
    Backend API        & ASP.NET Core 10, C\#, MediatR, EF Core, Serilog \\
    Database           & PostgreSQL 16 with TimescaleDB extension \\
    Messaging          & Eclipse Mosquitto (MQTT), SignalR + Redis \\
    Deployment         & Docker, Docker Compose, Nginx (production) \\
    \bottomrule
  \end{tabular}
\end{table}

\section{Smartwatch Application}

The Wear OS application (\texttt{watchapp/}) targets Samsung Galaxy Watch 8. Key modules:

\begin{itemize}[leftmargin=*]
  \item \texttt{VitalsSensorCoordinator} --- manages Samsung SDK sensor sessions for HR, SpO$_2$, skin temperature.
  \item \texttt{MotionSensorHub} --- step counter and fall detection via \texttt{TYPE\_ACCELEROMETER}.
  \item \texttt{VitalsMonitorService} --- Android foreground service keeping sensors and MQTT alive in background.
  \item \texttt{MqttManager} --- publishes JSON payloads to \texttt{vitals/\{patientId\}/data}.
  \item \texttt{VitalsMonitorScreen} --- Compose UI with sensor tabs (HR / Temp / SpO$_2$) and Start/Stop controls.
\end{itemize}

\projectfigure{watch-heart-rate.png}{6.5cm}{Wear OS application measuring heart rate in real time}{fig:watch-hr}

\projectfigure{watch-spo2-temp.png}{6.5cm}{Wear OS application showing SpO$_2$ or temperature measurement}{fig:watch-spo2}

\subsection{MQTT Payload Format}

Each published message contains:

\begin{verbatim}
{
  "patientId": "uuid",
  "deviceId": "uuid",
  "heartRateBpm": 92.0,
  "spO2Percent": 97.5,
  "temperatureC": 36.8,
  "skinTemperatureC": 33.2,
  "stepsCount": 1240,
  "caloriesBurned": 49.6,
  "fallDetected": false,
  "isWearing": true
}
\end{verbatim}

\section{Backend Implementation}

\subsection{MQTT Ingestion}

\texttt{MqttBackgroundService} connects to Mosquitto, subscribes to \texttt{vitals/+/data}, and forwards each message to \texttt{IngestVitalCommand} via MediatR. Patient and device IDs are validated as GUIDs before persistence.

\subsection{Alert Evaluation}

On each \texttt{VitalRecordedEvent}, the handler loads the patient's \texttt{AlertThreshold} and compares:

\begin{itemize}[leftmargin=*]
  \item Heart rate vs.\ min/max bounds
  \item SpO$_2$ vs.\ minimum threshold
  \item Systolic BP vs.\ maximum
  \item Temperature vs.\ maximum
  \item \texttt{fallDetected} flag (critical alert)
\end{itemize}

Created alerts raise \texttt{AlertTriggeredEvent}, which triggers FCM push to assigned doctors/relatives and SignalR broadcast.

\subsection{Real-Time Hub}

\texttt{VitalsHub} exposes two client callbacks: \texttt{ReceiveVitals} and \texttt{ReceiveAlert}. Clients connect with a JWT access token and subscribe to a specific patient's stream.

\section{Android Mobile Application}

The Android app (\texttt{android/}) provides:

\begin{itemize}[leftmargin=*]
  \item User authentication (login/register) with JWT token storage.
  \item Patient list for doctors; linked-patient view for relatives.
  \item Live vitals display via \texttt{VitalsSignalRClient}.
  \item Push notification handling via \texttt{RpmFirebaseMessagingService}.
  \item Chat functionality between care-team members (\texttt{ChatHub}).
\end{itemize}

\projectfigure{android-login.png}{7cm}{Android mobile application login screen}{fig:android-login}

\projectfigure{android-vitals.png}{7cm}{Android application displaying patient vitals or patient list}{fig:android-vitals}

\section{Web Application}

The React web app (\texttt{web-app/}) includes:

\begin{itemize}[leftmargin=*]
  \item Landing, login, and registration pages.
  \item Role-based dashboard with patient overview and admin panel.
  \item \texttt{PatientMonitor} page --- live SignalR vitals grid (HR, SpO$_2$, BP, temperature, fall status, watch wearing status).
  \item Heart rate dedicated monitor view.
\end{itemize}

\projectfigure{web-dashboard.png}{7cm}{Web dashboard showing multiple vital sign tiles and patient overview}{fig:web-dashboard}

\projectfigure{web-patient-monitor.png}{7cm}{Web patient monitor with real-time SignalR vital updates}{fig:web-monitor}

\section{Security}

\begin{itemize}[leftmargin=*]
  \item JWT Bearer authentication on all protected API endpoints and SignalR hubs.
  \item Passwords hashed before storage; role-based authorization on controllers.
  \item MQTT broker configurable for TLS (port 8883) in production.
  \item HTTPS termination via Nginx reverse proxy in production deployment.
\end{itemize}

% ============================================================
\chapter{Testing and Results}
\label{ch:testing}
% ============================================================

\section{Testing Methodology}

Testing was conducted at three levels:

\begin{enumerate}[leftmargin=*]
  \item \textbf{Unit/Integration Testing:} Backend command handlers, MQTT payload parsing, alert threshold logic.
  \item \textbf{Component Testing:} Watch sensor readings, MQTT publish/subscribe, SignalR client connections.
  \item \textbf{End-to-End Testing:} Full pipeline from watch on wrist to dashboard update and alert notification.
\end{enumerate}

\section{Test Environment}

\begin{itemize}[leftmargin=*]
  \item Hardware: Samsung Galaxy Watch 8, Android phone, development PC.
  \item Software: Docker Compose stack (API, PostgreSQL, Redis, Mosquitto).
  \item Network: Local Wi-Fi with MQTT broker at \texttt{host:1883}.
\end{itemize}

\projectfigure{testing-setup.png}{6.5cm}{Physical testing setup with smartwatch, mobile device, and live dashboard}{fig:testing-setup}

\projectfigure{docker-services.png}{5.5cm}{Docker Compose services running (postgres, redis, mosquitto, API)}{fig:docker}

\section{Functional Test Results}

\begin{table}[H]
  \centering
  \caption{Summary of functional test cases}
  \label{tab:testresults}
  \begin{tabular}{@{}p{5.5cm}cc@{}}
    \toprule
    \textbf{Test Case} & \textbf{Expected} & \textbf{Result} \\
    \midrule
    Watch publishes HR via MQTT       & Server receives \& stores record  & Pass \\
    Dashboard shows live HR update    & Value within 2 s                  & Pass \\
    SpO$_2$ below threshold           & Critical alert generated          & Pass \\
    Heart rate above maximum          & High-severity alert               & Pass \\
    Simulated fall (sharp wrist motion)& Fall alert triggered             & Pass \\
    Doctor receives FCM notification  & Push on linked device             & Pass \\
    JWT login and role access         & Correct page access per role      & Pass \\
    Historical vitals query           & Records returned from DB          & Pass \\
    \bottomrule
  \end{tabular}
\end{table}

\projectfigure{alert-notification.png}{6cm}{Alert notification displayed on mobile device or web dashboard}{fig:alert}

\projectfigure{swagger-api.png}{6.5cm}{Swagger UI documenting the RPM REST API endpoints}{fig:swagger}

\section{Performance Observations}

Under local network conditions, MQTT-to-dashboard latency averaged under 1.5 seconds. The foreground service on the watch maintained stable sensor readings during 30-minute continuous monitoring sessions. PostgreSQL with TimescaleDB handled time-series vital inserts without noticeable degradation at prototype data volumes.

\section{Discussion}

The system successfully demonstrates a complete RPM pipeline integrating mechatronic sensing (wearable hardware), IoT communication (MQTT), and software engineering (web/mobile/cloud). Fall detection proved sensitive to sharp wrist movements; future work should add machine-learning classification to reduce false positives. SpO$_2$ and skin temperature readings depend on proper watch fit and Samsung SDK availability.

% ============================================================
\chapter{Conclusion and Future Work}
\label{ch:conclusion}
% ============================================================

\section{Conclusion}

This graduation project designed and implemented a \textbf{Remote Healthcare Monitoring System} that integrates wearable technology, IoT messaging, and cloud-based services into a unified healthcare platform. The system continuously monitors heart rate, SpO$_2$, temperature, activity, and fall events from a Samsung Galaxy Watch, transmits data via MQTT to an ASP.NET Core backend, and presents real-time information to doctors, patients, and family members through web and mobile applications.

Automatic alert generation and push notification ensure that abnormal conditions are communicated promptly. The Docker-based deployment architecture provides a reproducible foundation for further development. The project demonstrates how mechatronics engineering principles---sensor integration, signal processing, and embedded systems---combine with modern software development to address real healthcare challenges.

\section{Future Work}

\begin{enumerate}[leftmargin=*]
  \item Integrate clinical-grade blood pressure and ECG sensors where hardware permits.
  \item Apply machine learning for improved fall detection and arrhythmia classification.
  \item Add FHIR-compliant data export for hospital EHR integration.
  \item Implement MQTT TLS and full HIPAA/GDPR compliance for production use.
  \item Conduct clinical trials to validate measurement accuracy against reference devices.
  \item Support additional wearable platforms (Apple Watch, Fitbit) via adapter layer.
  \item Add AI-based trend analysis and predictive alerts for early deterioration detection.
\end{enumerate}

% ============================================================
% REFERENCES
% ============================================================
\begin{thebibliography}{99}

\bibitem{rpm_survey}
O. Ajayi, A. Adewumi, and S. Misra, ``Remote Patient Monitoring: A Systematic Review of IoT-Based Healthcare Solutions,'' \textit{IEEE Access}, vol. 9, pp. 152852--152871, 2021.

\bibitem{rpm_chronic}
K. K. Venkatesh et al., ``Remote Monitoring of Chronic Disease Patients: A Meta-Analysis,'' \textit{Journal of Medical Internet Research}, vol. 22, no. 8, 2020.

\bibitem{fall_detection}
A. K. Bourke and G. M. Lyons, ``A Threshold-Based Fall-Detection Algorithm Using a Bi-Axial Gyroscope Sensor,'' \textit{Medical Engineering \& Physics}, vol. 29, no. 3, pp. 349--356, 2007.

\bibitem{mqtt_iot}
E. Al-Masri, K. Morley, and S. K. De, ``Real-Time IoT Stream Processing and Large-Scale Data Analytics for Smart City Applications,'' \textit{IEEE Internet of Things Journal}, vol. 4, no. 3, pp. 834--840, 2017.

\bibitem{samsung_sdk}
Samsung Electronics, ``Samsung Health Sensor SDK Developer Guide,'' 2024. [Online]. Available: \url{https://developer.samsung.com/health/sensor-sdk}

\bibitem{signalr}
Microsoft Corporation, ``ASP.NET Core SignalR Documentation,'' 2025. [Online]. Available: \url{https://learn.microsoft.com/en-us/aspnet/core/signalr/}

\bibitem{mosquitto}
Eclipse Foundation, ``Eclipse Mosquitto MQTT Broker Documentation,'' 2025. [Online]. Available: \url{https://mosquitto.org/documentation/}

\end{thebibliography}

% ============================================================
\appendix
\chapter{Source Code Repository Structure}
% ============================================================

\begin{verbatim}
Remote-Patiant-Monitoring/
|-- watchapp/          # Wear OS smartwatch application (Kotlin)
|-- android/           # Android mobile application (Kotlin)
|-- web-app/           # React web dashboard
|-- backend/
|   |-- RPM.Domain/        # Entities, enums, domain events
|   |-- RPM.Application/   # CQRS handlers, DTOs
|   |-- RPM.Infrastructure/# EF Core, MQTT, FCM, Redis
|   `-- RPM.API/           # REST API, SignalR hubs
|-- docker/            # Docker Compose infrastructure
`-- docs/
    `-- graduation-report/   # This report
\end{verbatim}

\chapter{API Endpoints (Summary)}

\begin{table}[H]
  \centering
  \caption{Selected REST API endpoints}
  \begin{tabular}{@{}llp{6cm}@{}}
    \toprule
    \textbf{Method} & \textbf{Endpoint} & \textbf{Description} \\
    \midrule
    POST & /api/auth/register    & Register new user \\
    POST & /api/auth/login       & Authenticate and receive JWT \\
    GET  & /api/vitals/\{patientId\} & Get vital history \\
    POST & /api/vitals/ingest    & Manual vital ingestion (testing) \\
    GET  & /api/alerts/\{patientId\} & List patient alerts \\
    PUT  & /api/alerts/threshold & Update alert thresholds \\
    GET  & /health               & Service health check \\
    \bottomrule
  \end{tabular}
\end{table}

\chapter{Default Alert Threshold Values}

\begin{table}[H]
  \centering
  \caption{Default per-patient alert thresholds (configurable)}
  \begin{tabular}{@{}ll@{}}
    \toprule
    \textbf{Parameter} & \textbf{Default Value} \\
    \midrule
    Minimum Heart Rate   & 40 bpm \\
    Maximum Heart Rate   & 130 bpm \\
    Minimum SpO$_2$      & 90\% \\
    Maximum Systolic BP  & 160 mmHg \\
    Maximum Diastolic BP & 100 mmHg \\
    Maximum Temperature  & 38.5\textdegree C \\
    \bottomrule
  \end{tabular}
\end{table}

\end{document}
